\section{Introduction}
\label{sec:introduction}

Antisemitism --- hostility, prejudice, or discrimination directed against Jewish people --- has a long and devastating history, and its contemporary manifestations on social media platforms present urgent challenges for automated content moderation~\citep{becker2025decoding}. The rise of online antisemitism has been documented across multiple studies, with sharp increases observed during geopolitical flashpoints such as the 2021 Israel--Gaza conflict and the period following October 2023~\citep{jikeli2023antisemitic}. Social media platforms serve as vectors for both overt hate speech (slurs, explicit calls for violence) and subtler, more insidious forms of anti-Jewish prejudice that cloak themselves in political commentary, conspiracy theories, or ironic detachment.

Automated detection of antisemitism is critically needed at scale, yet it poses challenges that distinguish it from general hate speech or toxicity classification in at least three fundamental ways.

\paragraph{The criticism--hate boundary.}
The International Holocaust Remembrance Alliance (IHRA) Working Definition of Antisemitism~\citep{ihra2016working}, which guides the annotation of our dataset, explicitly acknowledges that ``criticism of Israel similar to that leveled against any other country cannot be regarded as antisemitic.'' This creates a detection challenge with no parallel in general toxicity tasks: the model must distinguish between, for example, \emph{``Israel's settlement expansion violates international law''} (legitimate political criticism) and \emph{``Jews control American foreign policy to serve Israel''} (antisemitic conspiracy theory). The distinction hinges on whether the statement targets a \emph{state's policies} or ascribes negative traits to \emph{Jewish people as a collective} --- a nuance that requires genuine contextual understanding rather than keyword matching.

\paragraph{Coded language and evolving tropes.}
Antisemitic discourse has always been characterized by euphemism and coded reference. Classical tropes --- dual loyalty, world domination conspiracies, blood libel --- reappear in modern guise. The term ``ZioNazi,'' which trivializes the Holocaust by equating Zionism with Nazism, appears in 529 tweets in our dataset, with 88.3\% labeled antisemitic. Meanwhile, the slur ``kike'' appears in 283 tweets, but only 34.3\% of these are antisemitic --- the remainder being counterspeech that quotes the slur to denounce it, or meta-commentary about the word itself. These patterns illustrate that the mere presence of offensive terms is an unreliable predictor of antisemitic intent.

\paragraph{Class imbalance with heterogeneous sub-populations.}
In the GoldStandard2024 dataset~\citep{jikeli2024goldstandard}, only 1{,}953 of 11{,}311 tweets (17.3\%) are labeled antisemitic, yielding a 4.79:1 imbalance ratio. Crucially, the four keyword sub-populations have very different base rates --- Jews 11.4\%, Israel 16.1\%, Kikes 34.3\%, ZioNazi 88.3\% --- so a single overall accuracy or calibration number can mask large per-group disparities. This makes calibration and fairness genuinely separate problems from accuracy on the positive class.

\paragraph{Three-axis framing.}
The disjoint nature of these challenges leads us, contra most prior work, to study model behaviour jointly along three connected axes: (i)~\emph{accuracy} on the positive class, measured by positive-class F\textsubscript{1}; (ii)~\emph{calibration}, evaluated \emph{per sub-group} (group-conditional Expected Calibration Error) rather than only on average, since a model can be globally well-calibrated and severely miscalibrated on the smaller groups; and (iii)~\emph{counterfactual stability}, measured by the Counterfactual Flip Rate (CFR) under controlled identity-token swaps within a religiously-symmetric vocabulary (e.g.\ ``Jewish'' $\to$ ``Muslim/Christian/Hindu/Buddhist''). The three axes are not redundant: post-hoc Temperature Scaling~\citep{guo2017calibration}, the canonical calibration fix, provably cannot reduce the hard counterfactual flip rate at the decision threshold (\S\ref{sec:methodology}); and post-hoc multicalibration~\citep{hebertjohnson2018multicalibration}, while improving subgroup-conditional calibration, can degrade counterfactual invariance (\S\ref{sec:results}).

\paragraph{Approach.}
We propose \textbf{Counterfactual Calibration Invariance (CCI v2)}, a training-time loss that penalises the Jensen-Shannon divergence between predictions on a sentence and on its identity-token-swapped counterfactual, modulated by a per-group temperature head. CCI v2 generalises Counterfactual Logit Pairing~\citep{garg2019counterfactual,davani2021improving} along two axes: divergence in \emph{probability} space (rather than L\textsubscript{p} in logit space) and a per-group learnable temperature that lets the model modulate confidence locally. We pair the loss with a heuristic symmetry filter that rejects swaps involving slurs (where the swap fundamentally changes the meaning) or geopolitical proxies without a religiously-symmetric counterpart, ensuring that CCI v2 only enforces invariance on swaps that actually preserve the label. To position CCI v2 against the dominant alternative paradigm, we additionally benchmark prompted large language models --- Mistral-7B-Instruct~\citep{jiang2023mistral}, Qwen-2.5-7B-Instruct~\citep{qwen25}, Llama-3.1-8B-Instruct~\citep{llama31}, and Llama-3.3-70B (via the Groq inference API) --- under zero-shot, few-shot, and Guided Chain-of-Thought prompting~\citep{patel2025evaluating}.

\paragraph{Contributions.}
\begin{enumerate}[leftmargin=*,itemsep=2pt]
    \item \textbf{The headline observation: multicalibration vs.\ counterfactual invariance are in tension on text classifiers.} Applying multicalibration~\citep{hebertjohnson2018multicalibration} post-hoc to a strong baseline reduces the hard counterfactual flip rate by 8\% but \emph{increases the L1-norm soft flip rate by 96\%} (from 0.032 to 0.063). The canonical post-hoc fairness toolkit, while improving subgroup-conditional calibration, can substantially \emph{degrade} counterfactual invariance --- a tension between two well-respected fairness desiderata that, to our reading, has not been quantitatively documented for text classifiers. The broader incompatibility between calibration and equal error-rate fairness is established~\citep{pleiss2017fairness}; ours is the empirical specialisation to text counterfactual pair-invariance.
    \item \textbf{CCI v2: a training-time counterfactual loss with architecture-and-family transfer.} We propose CCI v2, a Jensen-Shannon divergence on identity-token-swapped counterfactuals within a religiously-symmetric vocabulary. A short structural lemma (Lemma~\ref{thm:ts-cfr}) shows that scalar post-hoc Temperature Scaling cannot reduce the hard CFR at threshold $1/2$, so CFR reduction requires a training-time intervention (like CCI v2) or a non-monotone post-hoc method (like multicalibration, with the regression in contribution 1). Empirically, CCI v2 reduces hard CFR by 49\% on DeBERTa-v3-base~\citep{he2023debertav3} (184M), 55\% on DeBERTa-v3-large (440M, same family), and 73\% on RoBERTa-large~\citep{liu2019roberta} (355M, different family) versus a strong focal+masking baseline at iso-architecture, with positive-class F\textsubscript{1} preserved or improved at scale (0.72 best). All four base-comparator CFR rejections survive Holm-Bonferroni; \emph{within-architecture} transfer is significant on both large backbones for CFR\textsubscript{hard} and CFR\textsubscript{soft, L\textsubscript{1}} under a separate Holm-Bonferroni family.
    \item \textbf{An honest fairness audit.} Under the Dixon~\citep{dixon2018measuring} sum-of-deviations FNED, the pivot CCI~v2 (fixed T, JS) sits within seed noise of every base-architecture baseline ($\Delta = +0.007$ vs.\ Davani $\lambda=0.5$, $p_{\text{boot}} = 0.86$; $\Delta = -0.004$ vs.\ Model~C, $p_{\text{boot}} = 1.0$). The cross-architecture FNED gaps ($+5.6\%$ on DeBERTa-large, $+7.3\%$ on RoBERTa-large) are also not Bonferroni-significant. The originally pre-registered CCI~v2 (learned T, JS) cell sits about $+5\%$ higher and is kept as an ablation comparator; the headline conclusion is that under the corrected aggregation, CCI~v2 trades nothing on FNED for the 49--73\% CFR reduction it gains.
    \item \textbf{Bimodal ECE under focal + F\textsubscript{1} early stopping (supporting).} Across six seeds, focal-loss training with F\textsubscript{1}-based early stopping yields a bimodal ECE distribution (4/6 seeds well-calibrated $\sim 0.05$, 2/6 poorly-calibrated $\sim 0.12$). We trace the bimodality to an interaction between the F\textsubscript{1}-peak epoch and the calibration trajectory, consistent with \citet{mukhoti2020calibrating} and \citet{tao2023pcs}. We report this as a supporting observation rather than a headline contribution.
    \item \textbf{Attribution mechanism check.} Integrated Gradients over six seeds shows that both Model~C and CCI~v2 reduce attribution mass on identity tokens relative to vanilla DeBERTa-v3. The additional CCI~v2 reduction over Model~C is statistically detectable but very small ($d_z=0.112$), locating its larger CFR gains primarily in pair-level counterfactual constraints rather than in a large further suppression of identity-token salience.
\end{enumerate}

The prompted-LLM comparison (Mistral-7B, Qwen-2.5-7B, Llama-3.1-8B, Llama-3.3-70B via Groq) and zero-shot cross-dataset transfer to the Jewish-target subsets of HateXplain~\citep{mathew2021hatexplain} and ToxiGen~\citep{hartvigsen2022toxigen} are part of a parallel evaluation track wired into the same pipeline; we report them in a companion artefact as they complete rather than as core contributions of the present paper.

\paragraph{Outline.}
Section~\ref{sec:related} reviews related work in hate-speech detection, counterfactual fairness, and antisemitism-specific NLP. Section~\ref{sec:methodology} describes the dataset, the three-axis evaluation protocol, the CCI v2 loss with its temperature head and symmetry filter, and the structural CFR-vs-TS argument. Section~\ref{sec:experiments} details implementation specifics and the multi-architecture sweep. Section~\ref{sec:results} presents the main results, the multicalibration-tension finding, and the Bonferroni-corrected significance suite. Section~\ref{sec:discussion} provides error analysis, the FNED Pareto trade-off discussion, and limitations. Section~\ref{sec:ethics} addresses ethical considerations including the IHRA-definition framing. Section~\ref{sec:conclusion} concludes.
